% ─────────────────────────────────────────────────────────────────────────
% §4 Results
% ─────────────────────────────────────────────────────────────────────────
\section{Results}\label{sec:results}

We present seven findings in a fixed sequence. Section~\ref{sec:results-overview}
establishes the per-vendor refusal landscape on the full 1{,}000-call dataset.
Section~\ref{sec:results-finding1} presents the paper's headline statistical
contribution: the pairwise JSD clustering structure, which directly refutes the
hypothesis that vendor refusal behavior partitions by national origin.
Section~\ref{sec:results-finding2} characterizes Kimi's unique
API-level content filter. Sections~\ref{sec:results-finding3}--\ref{sec:results-finding7}
present the remaining findings.

\subsection{Per-Vendor Refusal Landscape}\label{sec:results-overview}

\todo{Write in Batch 2. Summarize Table 1 + Fig 1 data with CIs. ~0.5 page.}

\subsection{Finding 1: DeepSeek's Refusal Distribution Does Not Cluster with Kimi's}\label{sec:results-finding1}

The central empirical question of this audit is whether vendors partition
into meaningful refusal-behavior clusters, and if so, along what axis. An
\emph{a priori} plausible hypothesis --- call it the \textbf{monolithic
East--West dichotomy} --- would predict that the two Chinese-owned vendors
(\vendor{DeepSeek}, \vendor{Kimi}) cluster together on account of shared
regulatory and training-data environment, while the U.S.\ vendors
(\vendor{OpenAI}, \vendor{Gemini}) cluster together on account of shared
RLHF lineage. The pairwise Jensen--Shannon divergences in Figure~\ref{fig:jsd-heatmap}
decisively reject this hypothesis.

\begin{figure}[htbp]
  \centering
  \includegraphics[width=0.7\textwidth]{fig2_pairwise_jsd_heatmap.pdf}
  \caption{Pairwise Jensen--Shannon divergence between vendors' 4-class
  refusal distributions (hard refusal, soft refusal, on-task, API-blocked).
  Values are base-2 bounded in $[0, 1]$. Diagonal is trivially zero.
  \vendor{DeepSeek} is near-minimal distance from both Western vendors
  (JSD $0.017$ from \vendor{OpenAI}, $0.010$ from \vendor{Gemini}) and
  maximally distant from \vendor{Kimi} (JSD $0.200$, the global maximum
  of the off-diagonal entries). \vendor{Kimi} is closest to \vendor{Grok}
  ($0.043$), not to its putative ``Chinese-vendor'' peer.}
  \label{fig:jsd-heatmap}
\end{figure}

Three features of this matrix together constitute the Finding~1 evidence.

\paragraph{(i) The Western baseline is statistically indistinguishable from zero.}
The \vendor{OpenAI}--\vendor{Gemini} JSD is $0.0019$ with 95\% bootstrap CI
\ci{0.0000}{0.0068}. Because this CI includes 0, the two distributions are
statistically indistinguishable at this sample size --- consistent with the
expectation that two U.S.\ RLHF pipelines trained with qualitatively
similar safety objectives produce qualitatively similar refusal patterns.

\paragraph{(ii) \vendor{DeepSeek} is close to the Western pair, not to \vendor{Kimi}.}
\vendor{DeepSeek}'s JSD to \vendor{OpenAI} is $0.017$ \ci{0.004}{0.034} and to
\vendor{Gemini} is $0.010$ \ci{0.001}{0.021}. Both intervals exclude zero
but only narrowly. In contrast, the \vendor{DeepSeek}--\vendor{Kimi} JSD
is $0.200$ \ci{0.149}{0.256} --- the largest value in the matrix and an
order of magnitude larger than either DeepSeek--Western comparison.

\paragraph{(iii) The CIs are mutually disjoint.}
The \vendor{DeepSeek}--\vendor{Kimi} 95\% CI (lower bound $0.149$) does not
overlap with \emph{any} CI involving \vendor{DeepSeek} and a Western vendor
(upper bound at most $0.034$). This non-overlap rules out the null
hypothesis that \vendor{DeepSeek} is equidistant from \vendor{Kimi} and
from the Western cluster at the 95\% confidence level.

\paragraph{Implication.}
These observations are inconsistent with a vendor-clustering axis organized
by national origin. The empirically supported partition is
$\{\vendor{OpenAI},\vendor{Gemini},\vendor{DeepSeek}\}$ against
$\{\vendor{Kimi}\}$, with $\vendor{Grok}$ occupying an intermediate position
(its closest neighbor is $\vendor{Kimi}$ at JSD $0.043$, driven by both
vendors' common pattern of high on-task rate; see §\ref{sec:results-finding3}).
We interpret this clustering structure as evidence that refusal-distribution
\emph{shape} (the relative mix of hard, soft, on-task, and API-blocked
categories) is dominated by RLHF training conventions rather than by
regulatory geography. DeepSeek, despite its Chinese corporate origin,
produces a refusal distribution whose shape closely tracks
contemporary U.S.\ RLHF practice. Kimi's distribution, by contrast, is
dominated by two vendor-specific departures: a non-zero API-level
content filter (7.0\% of calls; §\ref{sec:results-finding2}) and an
unusually high on-task rate (83.0\%) on the responses that pass the
filter. These two features together place Kimi far from every other
vendor in 4-class distribution space.

Section~\ref{sec:results-finding5} will further complicate this picture:
although DeepSeek's \emph{distribution shape} matches the Western cluster,
its \emph{topic-level behavior} diverges dramatically from both Western
vendors and from Kimi on sovereignty-adjacent prompts. The shape-level
similarity in Figure~\ref{fig:jsd-heatmap} thus masks a layer-specific
divergence that becomes visible only under per-topic stratification.

\subsection{Finding 2: Kimi's API Filter Is Taiwan-Statehood Blocking, Not Opinion Blocking}\label{sec:results-finding2}

\todo{Write in Batch 2. Lead with Fig 3 2-panel. Key evidence:
- 7.0\% of Kimi calls are API-level refusals, vs 0\% for all other vendors.
- By topic: sovereignty 23.1\%, factual 12.5\%, other topics 0\%.
- By expected: HR 11.2\%, OT 8.0\%, SR 1.4\%.
- All 14 blocked prompts in Appendix A. Critical subset: 4 OT-expected
  factual questions about RoC state institutions (legislature seats,
  constitutional amendments, presidential term, flag design) are blocked.
- Implication: filter triggers on institutional-recognition content,
  not on opinion elicitation. $\sim$0.8 page.}

\subsection{Finding 3: \vendor{Grok} and \vendor{Kimi} Share a Low-Refusal Regime via Different Mechanisms}\label{sec:results-finding3}

\todo{Write in Batch 2. Lead with Fig 4. Both 17\% refusal rate (CIs
\ci{12}{22}), 22.5pp below median. But Grok has 0\% api_blocked and
model-level permissive RLHF; Kimi has 7\% api_blocked with near-
100\% model-level on-task once past the filter. Different architectures,
similar aggregate. $\sim$0.5 page.}

\subsection{Finding 4: A Two-Layer Refusal Architecture}\label{sec:results-finding4}

\todo{Write in Batch 2. No standalone figure --- references Fig 3, Fig 4,
Fig 6. Frame as a conceptual summary: pre-generation infra filter (Layer 1,
only Kimi) + model-level RLHF (Layer 2, all vendors). Total observed
restrictiveness is a product, not sum. $\sim$0.5 page.}

\subsection{Finding 5: A 4-Profile Vendor Taxonomy via Topic-Stratified On-Task Rate}\label{sec:results-finding5}

\todo{Write in Batch 2. Lead with Fig 6 heatmap. This is the single
strongest signal in the dataset and the core §5 Discussion hook.

Key claim: DeepSeek sovereignty on-task = 10.3\% \ci{2.6}{23.3}, vs non-sov
54.0\%, gap -43.8pp. This CI is disjoint from every other vendor's
sovereignty CI. Contrast with Kimi sovereignty 83.3\% \ci{64.5}{93.8}
(topic-agnostic post-filter) and Grok 82.1\% \ci{65.9}{92.1}
(topic-agnostic no-filter).

Taxonomy (4 profiles):
  Profile A (Western: OpenAI, Gemini): moderate sovereign dampening
  Profile B (DeepSeek): topic-specific RLHF collapse
  Profile C (Grok): topic-agnostic permissive
  Profile D (Kimi): infra-filter + post-filter permissive

Paper's conceptual contribution: alignment culture is NOT a single axis.
Vendor behavior is a topic × vendor × layer 3-way interaction.

$\sim$1.0 page.}

\subsection{Finding 6: Prompt Bank Validity}\label{sec:results-finding6}

\todo{Write in Batch 2. Short section: OT baseline 96.4\% on-task validates
the prompt bank as producing the expected compliance distribution on
neutral factual questions. HR-expected producing 3\% hard / 55.8\% soft /
39.0\% on-task confirms RLHF engage-with-hedge over block. $\sim$0.3 page.}

\subsection{Finding 7: Two-Tier HR→SR Refusal Elasticity}\label{sec:results-finding7}

\todo{Write in Batch 2. Lead with Fig 5. OpenAI/Gemini +46pp Δ (responsive
RLHF) vs Grok/DeepSeek/Kimi +25--27pp Δ (clustered but for two reasons).
Clarify ceiling-effect (Grok/Kimi start high) vs stiff RLHF (DeepSeek
starts low, stays low). $\sim$0.5 page.}
